\subsection*{寻找没有偏离原直线的方向}
一般向量经过矩阵后方向与长度都变，但某些方向只伸缩，仍在原直线上。找到这些方向，就可能把复杂动作拆成几种互不干扰的伸缩。特征向量是这些方向上的非零向量；倍率可以为负，表示反向，也可以为零，表示该方向被压掉。

对 $A=\left[\begin{smallmatrix}2&1\\1&2\end{smallmatrix}\right]$，$(1,1)$ 变为 $(3,3)$，$(1,-1)$ 不动。任意输入都可用这两方向拼成，每作用一次只需把第一份系数乘以 $3$。多次作用后的主要增长方向因此一眼可见。

\subsection*{特征方程为什么出现}
要求 $Av=\lambda v$ 且 $v\ne0$，等价于要求 $A-\lambda I$ 有非零核。也就是从原动作中减去统一伸缩之后，能否压掉某个方向。对方阵，这要求行列式为零，所以保留方向的问题转成关于倍率的多项式方程。

求出根后仍需找向量。四分之一圈的实平面旋转没有实的保留直线，所以没有实特征向量；非平凡剪切可能只有一条保留直线，不足以组成基。能对角化要求足够多独立特征向量，不能只检查特征值存在。

在矩阵实验中选择对称示例，转动输入方向，预测哪两条线被保留。再把任意输入拆成这两个方向，预测四次作用的结果。计算应来自系数的幂，而非逐次展开四个矩阵；这是几何理解带来的实际收益。
